\subsection{Segmentation-First Routing to Two Hough Searches}
\label{sec:segmentation-guided-hough}

The preceding watershed experiment depended on existing Hough detections to
locate crowded windows and on distance-transform maxima to propose missing
markers. Its failure to recover the striped blue ball above the black ball
in Table4 motivated a different order of operations. The next experiment
constructed foreground components before running Hough, identified potential
clusters from their area and shape, and used different Hough settings inside
and outside those windows. A separate watershed check then assessed whether
the newly generated markers improved candidate validation.

Three methods were compared: the current standard Hough detector;
segmentation-guided Hough with circular-arc validation; and the same guided
candidates followed by watershed-basin validation. All methods used the same
known table polygon and the same final pocket filter. Detection proposals
were generated before loading ground-truth annotations. No annotated ball
centre was used to select a component, seed a marker, or accept a circle.
The main application was not changed.

\subsubsection{Foreground components before any Hough call}

Foreground was estimated over the table using $128\times128$ image tiles.
Within each tile, every second pixel in each direction was eligible as a
cloth sample when its HSV values satisfied $35\leq H\leq95$,
$S\geq35$ and $V\geq30$. At least 30 samples were required. The
component-wise median of their encoded 8-bit Lab values supplied the local
cloth colour. The foreground-distance threshold was the median sample
distance plus 2.5 times its median absolute deviation, clipped to $[16,45]$,
as in the preceding experiment. A global cloth estimate provided a fallback
for tiles without enough samples. If no global model was available, no
cluster was declared and standard Hough was retained everywhere.

No Hough-based exclusion disks or foreground cores were used at this stage.
This distinguishes the experiment from the earlier Hough-first watershed
prototype. Local cloth estimation can still confuse a green ball with the
cloth, but an isolated ball omitted from the foreground is not automatically
deleted: standard Hough still examines the original table-masked image.

The foreground construction inset the table polygon with a $41\times41$
elliptical erosion to suppress rail and shadow connections. A $3\times3$
elliptical closing filled small gaps, enclosed holes were filled, and the
mask was restricted to the inset table. Components smaller than 80 pixels
were discarded. An initial engineering run with an 11-pixel inset had
connected foreground along rails into nearly table-wide search windows.
Its outputs were retained separately; the reported configuration uses the
41-pixel inset. This adjustment was based on inspecting the mask failure,
but the experiment remains exploratory and should not be presented as an
untouched held-out evaluation.

An image-coordinate size guide was constructed without Hough detections:
the expected radius increased linearly from 15 to 30 pixels between the
minimum and maximum source-image y coordinates of the table corners.
This used the existing radius bounds and camera geometry, not annotated
ball sizes. It is a heuristic; actual apparent radii need not follow
that exact linear relation.

A foreground component was considered a potential cluster when its area
was at least 1.6 times the expected disk area, or when its bounding-box
aspect ratio was at least 1.45 and its area was at least 0.8 expected disk
areas. Components covering over 10\% of the table, or long rail-like
components with aspect ratio above four and more than 60\% of pixels in
the table's boundary band, were excluded. A window extended each selected
component's bounding box by 1.25 expected radii. Overlapping windows were
merged. The final run produced 32 windows across seven images; some
represented shadows or single balls rather than actual multi-ball clusters.

\subsubsection{Routing standard and permissive circle searches}

Standard Hough retained the original settings: $dp=1.5$, minimum centre
distance 18 pixels, upper Canny threshold 75, accumulator threshold 35,
and radius bounds 15--30 pixels. It was computed on the unmodified
table-masked image and then routed by candidate centre. This avoids
introducing artificial image edges by painting cluster interiors black
before circle detection.

Standard candidates centred in the interior of a cluster bounding box
were replaced by results from its local permissive search. An overlap
margin of half an expected radius retained standard candidates near
the region boundary, while the padded local crop covered the same area.
Thus routing did not intentionally leave a gap between the two searches.

The local search retained $dp=1.5$, Canny threshold 75 and the original
$5\times5$, $\sigma=1.5$ preprocessing, but reduced the accumulator
threshold to 22 and minimum centre distance to 12 pixels. Its radius
bounds were $\max(12,\lfloor0.7r_e\rfloor)$ and
$\min(32,\lceil1.6r_e\rceil)$, where $r_e$ is the window's size guide.
This widens the plausible range modestly instead of treating occlusion
as evidence that the physical ball radius has become much smaller.
The exact bounds for each window were saved to CSV.

Local searches used the original colour crops, not binary foreground
images. Foreground was a routing and support cue; Hough still measured
image edges. Each proposal was checked using 72 samples around its
predicted circumference. A sample was supported when a Canny edge was
within 2.5 pixels and the absolute cosine between its Sobel gradient
direction and the circle's radial direction was at least 0.6. This
reduces support from edges crossing a proposed circle tangentially or
from unrelated internal markings.

At least 30\% of circumference samples had to be supported. A contiguous
supported arc of at least 75 degrees was also required, allowing isolated
one-sample gaps to be bridged for the arc-length calculation only.
The raw support percentage was not increased by that bridging. At least
35\% of the proposed disk had to overlap foreground, and the centre
had to lie inside the table. These tests allow a partly occluded circle
to lack edge evidence along its hidden side.

Candidates passing these checks were ordered by
$0.7\times\text{edge support}+0.3\times\text{arc angle}/360$.
A candidate was considered a duplicate only when its centre was closer
than 0.8 times the smaller radius to an already selected circle.
This replaces the earlier distance-peak experiment's much wider
1.15-larger-radius exclusion, which could suppress occluded neighbours.
Standard detections retained by routing were selected first. If a
window accepted no local proposal, available standard circles in that
window were retained as an explicit empty-search fallback. This decision
did not consult ground truth.

\subsubsection{Separate watershed ablation}

The optional watershed stage used accepted circle centres as distinct
markers. This directly addresses the earlier missing-marker limitation:
the new markers came from local image-edge evidence, not from maxima
of the merged foreground silhouette. Small cores of half a detected
radius were added to the watershed foreground around these accepted
centres, after cluster selection had already finished. Foreground
markers had radius three pixels; background lay outside a $7\times7$
foreground dilation or outside the table. Watershed used the original
BGR crop smoothed with a $3\times3$ Gaussian, $\sigma=0.8$.

For each accepted local candidate, the resulting label area overlapping
the watershed foreground had to fall between 0.15 and 2.0 times its
expected disk area. This was a candidate-rejection check, not a new
centre-estimation method. Standard detections and explicit fallback
detections were not removed by this check. Centres and radii were not
changed. Comparing the guided variant before and after this check
isolates this particular contribution of watershed; it does not test
every possible use of watershed contours or masks.

\subsubsection{Table4 recovery and overall results}

% Generated by experiments/writeRecognitionReport.py from saved CSVs.
\begin{table}[htbp]
\centering
\small
\caption{Table4: recovering ball 10 does not imply an improved overall detection result; other lost matches and false candidates must also be counted.}
\label{tab:segmentation-hough-table4}
\begin{tabular}{lrrrcr}
\toprule
Method & TP & FP & FN & Ball 10 recovered & Correct no. \\
\midrule
Standard Hough & 14 & 0 & 1 & No & 13 \\
Segmentation-guided Hough & 14 & 2 & 1 & Yes & 12 \\
Guided Hough + watershed & 14 & 2 & 1 & Yes & 12 \\
\bottomrule
\end{tabular}
\end{table}
\FloatBarrier


The guided search recovered the previously missed ball 10 in Table4.
Its detected source centre was $(1479.75,104.25)$ with radius 19.5 pixels.
After perspective transformation its centre error was 7.00 table pixels,
within the fixed 15-pixel matching tolerance. The watershed check retained
this candidate. Thus visible circular-edge evidence could recover the
ball that the earlier distance-peak marker strategy missed. The fixed
recognition pipeline correctly predicted striped for this candidate,
but returned an unknown number; recovery of its location did not imply
correct number recognition.

\begin{figure}[htbp]
\centering
\includegraphics[width=\linewidth]{segmentation_hough_table4.png}
\caption{Evaluation-only close-up centred on the missed Table4 annotation.
The standard detector misses ball 10, while the local permissive search
recovers it and supplies a separate watershed marker. Green circles
are matched detections, red circles are false detections and the blue
cross marks a missed annotation. Ground truth selected this display
crop after detection; it did not determine any proposal or marker.}
\label{fig:segmentation-hough-table4}
\end{figure}
\FloatBarrier

However, the same Table4 run lost the baseline match to ball 7 and
introduced additional false detections. Before watershed it produced
14 true and two false detections, compared with the baseline's 14 true
and zero false detections. The watershed area check retained the same
16 circles in Table4. This is a change in which balls were found,
not an overall Table4 detection improvement.

% Generated by experiments/writeRecognitionReport.py from saved CSVs.
\begin{table}[htbp]
\centering
\small
\caption{Segmentation-first detection comparison on Table1--6. Recoveries and losses are annotated-ball identities relative to standard Hough; correct-number counts use the same fixed K=6/CCA recognition pipeline for all methods.}
\label{tab:segmentation-hough-summary}
\begin{tabular}{lrrrrrrr}
\toprule
Method & TP & FP & FN & Precision & F1 & Recovered/lost & Correct no. \\
\midrule
Standard Hough & 86 & 2 & 4 & 97.73\% & 96.63\% & 0/0 & 75 \\
Segmentation-guided Hough & 86 & 16 & 4 & 84.31\% & 89.58\% & 3/3 & 71 \\
Guided Hough + watershed & 83 & 11 & 7 & 88.30\% & 90.22\% & 3/6 & 69 \\
\bottomrule
\end{tabular}
\end{table}
\FloatBarrier

% Generated by experiments/writeRecognitionReport.py from saved CSVs.
\begin{table}[htbp]
\centering
\small
\caption{Annotation identities recovered or lost by the segmentation-guided Hough variant before the optional watershed check.}
\label{tab:segmentation-hough-changes}
\begin{tabular}{lrl}
\toprule
Image & Ball & Change \\
\midrule
Table1.jpeg & 7 & recovered \\
Table1.jpeg & 9 & lost \\
Table3.jpeg & 10 & recovered \\
Table4.jpeg & 7 & lost \\
Table4.jpeg & 10 & recovered \\
Table6.jpeg & 10 & lost \\
\bottomrule
\end{tabular}
\end{table}
\FloatBarrier


Across Table1--Table6, guided Hough recovered three previously missed
annotations: ball 7 in Table1 and ball 10 in Table3 and Table4. It
also lost three previously matched annotations: ball 9 in Table1,
ball 7 in Table4 and ball 10 in Table6. Total true positives therefore
remained 86, with four misses. False detections increased from two
to 16, reducing precision from $97.73\%$ to $84.31\%$ and F1 from
$96.63\%$ to $89.58\%$. The mean matched-centre error also increased
from 4.91 to 5.29 table pixels.

Watershed reduced the guided variant's false detections from 16 to 11,
but removed three additional correct matches. It therefore produced
83 true positives, seven misses, precision $88.30\%$ and F1 $90.22\%$.
It was slightly better than unfiltered guided Hough in F1, but worse
than standard Hough. On AllBalls, guided Hough retained 16 true
detections while adding four false detections; watershed retained
14 true and three false detections, losing two reference balls.

% Generated by experiments/writeRecognitionReport.py from saved CSVs.
\begin{table}[htbp]
\centering
\small
\caption{Per-image true and false detections. AllBalls is the reference image, reported separately from the six comparison images.}
\label{tab:segmentation-hough-images}
\begin{tabular}{lrrrrrr}
\toprule
Image & Base TP & FP & Guided TP & FP & Guided+WS TP & FP \\
\midrule
AllBalls.jpeg & 16 & 0 & 16 & 4 & 14 & 3 \\
Table1.jpeg & 14 & 1 & 14 & 4 & 14 & 1 \\
Table2.jpeg & 16 & 1 & 16 & 5 & 14 & 4 \\
Table3.jpeg & 15 & 0 & 16 & 4 & 16 & 3 \\
Table4.jpeg & 14 & 0 & 14 & 2 & 14 & 2 \\
Table5.jpeg & 14 & 0 & 14 & 0 & 13 & 0 \\
Table6.jpeg & 13 & 0 & 12 & 1 & 12 & 1 \\
\bottomrule
\end{tabular}
\end{table}
\FloatBarrier


Recognition was frozen to configuration 2153 plus the selected solid-only
CCA correction, using ordinary circular masks. The baseline correctly
classified 82 types and 75 numbers among its 86 matched balls. Guided
Hough produced 78 correct types and 71 numbers, while guided Hough
with watershed produced 75 types and 69 numbers. Changed local centres
and radii can alter extraction and colour proportions even for an
annotation still matched by both detectors. These scores therefore
measure the downstream effect of the complete detection alternatives,
not only the identities of recovered balls. Recognition thresholds
were previously tuned on Table1--Table6, so these remain exploratory
comparison scores rather than independent recognition test estimates.

\subsubsection{Interpretation and reproducibility}

The experiment supports the narrow hypothesis that a more permissive
local edge search can supply a missing marker for the visible ball
10/ball 8 pair. It does not support replacing the current detector
with this complete prototype. Segmentation sometimes groups shadows
or single balls as clusters, local Hough accepts plausible circular
fragments from background and internal features, and changing an
existing crop can lose an otherwise acceptable annotation match.
The watershed area check also rejects some correct markers because
colour boundaries within a ball are not identical to object boundaries.
These costs must be counted alongside the visually convincing recovery.

A stronger follow-up could preserve confirmed standard centres in a
cluster and use local proposals primarily to supplement missing balls,
with stricter background and arc-consistency tests. That would be a
different comparison and was not silently substituted for the requested
two-route experiment. No segmentation IoU is reported because the
dataset supplies centres and numbers rather than pixel-level masks.

The command is:
\begin{quote}
\code{./build-local/BallExperiments --segmentation-hough}
\end{quote}
An optional output path preserves separate runs. Options are defined in:
\begin{quote}
\code{experiments/segmentationGuidedDetection.h}
\end{quote}
The default output directory is:
\begin{quote}
\code{output/experiments/segmentation_guided_detection/}
\end{quote}
It contains pooled/per-image scores, per-annotation coverage and
recovery/loss records, candidate predictions, full option values,
cluster windows and routing metadata, radial edge/arc evidence, and
watershed acceptance/rejection reasons. Diagnostic sheets show the
foreground before Hough, routing interiors, raw local proposals,
accepted local circles and watershed labels. Recovered annotations
receive additional close-up comparisons for evaluation only.

Synthetic checks confirmed that changing the local Hough threshold did
not change the foreground or routing mask, that accepted local candidates
met their radial-arc criteria, and that watershed output was a subset
of the guided circles. Empty local searches and absent cloth models
were checked for baseline fallback. All exported pooled/per-image scores
and annotation changes were reconciled with predictions, and the standard
control matched the preceding experiment's centres and recognition results.

Report tables and the Table4 figure are refreshed with:
\begin{quote}
\code{python3 experiments/writeSegmentationGuidedReport.py}
\end{quote}
The surrounding prose records this run's settings and results and
needs review after a rerun with different options. All implementation
changes remain in the experiment target; production detection and
recognition functions and their parameters are unchanged.
